% !TeX program = xelatex
% UTF-8; compile with XeLaTeX from this directory or the repository root.
% Source: Google Drive / notebook / anemia & hyperbilirubinemia.
% Images are unchanged originals, ordered by their original filenames.
% pic01.jpg = 20261016_110942358.jpg; Drive ID 16wYjcta35S2JoPWzm-7nN12H63B3d4tT.
% pic02.jpg = 20261016_110949909.jpg; Drive ID 1WXxo0aREa2GkQG0dVyy_C2afZmom0CFt.
% pic03.jpg = 20261016_110954282.jpg; Drive ID 1l7rPWGQ_jDQz6HJ9M45K7056O-m18eMG.
% pic04.jpg = 20261016_110959924.jpg; Drive ID 1AW1wmuGl43xtMfW845WR-Q1fuJAgCnB8.
% pic05.jpg = 20261016_111010474.jpg; Drive ID 1xxf1aVSj9GndHwTIgwOpCbzbRVBxBwaO.
% pic06.jpg = 20261016_111017215.jpg; Drive ID 1qPJZsSZ213X0bUqnmQCSqSucbmtpmCbx.
% pic07.jpg = 20261016_111021665.jpg; Drive ID 1_Zvo4gwEEN8i4okNpCspbErteNLzUEUI.
% pic08.jpg = 20261016_111026544.jpg; Drive ID 1Axs7xPCLYPq57wPCP58gYHQb6Anuc1eq.
% pic09.jpg = 20261016_111029988.jpg; Drive ID 1yQnBPiE8sTz-U3YV0iVRIyxeNLm7tJvy.
% pic10.jpg = 20261016_111034711.jpg; Drive ID 14gqjioHuZHFrYmHhi0juTWIqu7ckTeM5.
% pic11.jpg = 20261016_111040186.jpg; Drive ID 1rtSAtc7SwdXEObCQnKY3hWBBjfI0a0HI.
% pic12.jpg = 20261016_111045012.jpg; Drive ID 1RMmxkjXyHzXLaKYXhnuMkn9nf_R3TWeL.
% pic13.jpg = 20261016_111049397.jpg; Drive ID 1MZbaeeniALWq49IAvQsqK4ndH-VaHg8D.
% Define \AnemiaHyperbilirubinemiaInNotebook before input into an existing
% XePersian document (with amsmath, amssymb, and graphicx already loaded).
\ifdefined\AnemiaHyperbilirubinemiaInNotebook
  \let\AHfinish\relax
\else
  \def\AHfinish{\end{document}}
  \documentclass[12pt,a4paper]{article}
  \usepackage[margin=22mm]{geometry}
  \usepackage{amsmath,amssymb}
  \usepackage{graphicx}
  \usepackage{enumitem}
  \usepackage{xepersian}
  \IfFontExistsTF{Amiri}{\settextfont{Amiri}}{\settextfont{DejaVu Sans}}
  \setlatintextfont{DejaVu Serif}
  \setlength{\parindent}{0pt}
  \setlength{\parskip}{5pt}
  \setlength{\emergencystretch}{3em}
  \begin{document}
\fi
\providecommand{\AHmark}[1]{\textbf{#1}}
\providecommand{\AHunclear}[1]{\textbf{[خوانش نامطمئن: #1]}}
\providecommand{\AHmath}[1]{\lr{\(#1\)}}
\providecommand{\AHpage}[1]{\clearpage\section*{تصویر #1}}

\begin{center}
{\Large\bfseries ۶۲ ـ آنمی و هایپربیلی‌روبینمی}\par
\lr{Anemia \& hyperbilirubinemia}
\end{center}
\textbf{یادداشت رونویسی:}
متن فارسی و انگلیسی، اعداد، حاشیه‌نویسی‌ها و ارتباط شاخه‌های دست‌نویس به ترتیب
۱۳ تصویر بازنویسی شده‌اند. عبارت‌های هایلایت‌شده پررنگ‌اند و خوانش نامطمئن مشخص شده است.
این سند رونویسی یادداشت‌هاست؛ اعداد و عبارت‌های منبع بدون تصحیح علمی نقل شده‌اند.
اصل کامل هر ۱۳ تصویر در پیوست همین فایل آمده است.

سربرگ چاپی همهٔ صفحات: \lr{Subject:}، \lr{Year:}، \lr{Month:} و \lr{Date:}.
خانه‌های تاریخ خالی‌اند. سربرگ دست‌نویس تصویر ۱ «۶۲ ـ آنمی و هایپربیلی‌روبینمی»
و تصاویر ۲ تا ۱۳ «۶۲» است. اعداد چاپی کنار سطرها ۱ تا ۹ و شماره‌های دست‌نویس
پایین برگه‌ها به ترتیب ۱ تا ۱۳ هستند.

\AHpage{۱}
\subsection*{هماتوپوئز و هموگلوبین نوزاد}
\begin{itemize}
\item هماتوپوئز رویان از روز \AHunclear{۲…ام؛ انتهای عدد روز در دست‌نویس روشن نیست} آغاز می‌شود.
\item در \AHmark{سه‌ماههٔ سوم}، مغز استخوان تبدیل به محل اصلی هماتوپوئز می‌شود.
\item \lr{Hb} در هفتهٔ ۱۲ام بارداری: \lr{8--10 g/dL}؛
در هفتهٔ ۳۴ام بارداری: \AHmark{\lr{14.5--18 g/dL}}.
\item پس از تولد، \lr{Hb} تا ۶–۱۲ ساعت افزایش می‌یابد، سپس
\AHmark{تا ماه ۳–۶ تولد} به \AHmark{\lr{11--12 g/dL}} کاهش می‌یابد.
\item در نوزادان نارس، در ماه ۱–۲ به پایین‌ترین حد می‌رسد؛ زیر ۳۲ هفته.
\item \lr{RBC}های جنین و نوزاد: طول عمر کمتر، \AHmark{۷۰–۹۰ روز}؛
\lr{MCV} بالاتر، \AHmark{۱۱۰–۱۲۰}.
\item هنگام تولد، \AHunclear{۶۰–۹۰٪؛ رقم پایین بازه شبیه ۶۰ نوشته شده است}
هموگلوبین از نوع \lr{F} است؛ در ۴ماهگی: ۵٪.
\item حجم خون در نوزاد ترم: \lr{100 mL/kg}.
\item \AHmark{تأخیر در کلمپ بند ناف ریسک پلی‌سیتمی و زردی را افزایش می‌برد.}
در منبع، روی «پلی‌سیتمی» و «زردی» به ترتیب ۱ و ۲ نوشته شده است.
\end{itemize}
\subsection*{آنمی فیزیولوژیک}
\begin{itemize}
\item در \AHmark{ماه ۲–۳ تولد} شیرخوار ترم و ماه ۱–۲ تولد پره‌ترم.
\item علت آنمی فیزیولوژیک: افزایش اکسیژن بافتی، طول عمر کوتاه \lr{RBC}،
سطح پایین اریتروپوئتین.
\end{itemize}

\AHpage{۲}
\subsection*{علل آنمی نوزاد}
کاهش تولید \lr{RBC}، افزایش تخریب \lr{RBC}، از دست دادن خون.
\subsection*{کاهش تولید \lr{RBC}}
\begin{itemize}
\item علائم: رنگ‌پریدگی، رتیک پایین، فقدان هپاتواسپلنومگالی، اریتروئید هایپوپلازی \lr{BM}.
\item علل: سندرم نارسایی مغز استخوان:
\lr{Congenital RBC Aplasia (\AHmark{Blackfan-Diamond Anemia})}.
\item عفونت‌ها: روبلا، پارووویروس \AHmath{B_{19}}؛ سپسیس.
\item لوکمی همراه با تولد.
\end{itemize}
\subsection*{افزایش تخریب \lr{RBC}}
\begin{itemize}
\item همولیز ایمیون در رحم؛ منجر به \AHmark{اریتروبلاستوزیس فتالیس}.
\item همولیز ممکن است پس از تولد بروز کند.
\item تفاوت گروه خونی مادر و جنین در مادر حساس‌شده: همولیز.
\item \lr{IgG Ab} مادری علیه آنتی‌ژن‌های \lr{RBC} جنین.
\end{itemize}
\subsection*{چطور مادر حساس شده است؟}
\begin{itemize}
\item تزریق خون داشته؛ آمنیوسنتز.
\item سقط در سه‌ماههٔ اول.
\item \lr{EP}؛ جدا کردن دستی جفت.
\item مداخلات مختلف.
\item حتی در بارداری طبیعی هم ریسک حساسیت وجود دارد.
\end{itemize}
\begin{itemize}
\item \AHmark{ناسازگاری \lr{ABO} معمولاً بیش از یک آنمی خفیف ایجاد نمی‌کند.}
\item \AHmark{\lr{Hemolytic disease of newborn}} در موارد نادر:
آنمی شدید + هایپربیلی‌روبینمی.
\end{itemize}

\AHpage{۳}
\subsection*{همولیز ناشی از ناسازگاری گروه خون}
\begin{itemize}
\item آنتی‌بادی‌ای که بیشتر مادران با گروه خون \lr{O} علیه آنتی‌ژن \lr{A} و \lr{B} دارند:
\AHmark{واکنش حتی در بارداری اول هم ممکن است رخ دهد}؛ برخلاف \AHmark{\lr{Rh}}.
\item همولیز \lr{Rh} تحت تأثیر تعداد بارداری‌های قبلی: در بارداری اخیر شدیدتر تظاهر می‌کند.
\item \AHmark{همولیز \lr{Rh} شدیدتر از همولیز \lr{ABO} است.}
\item علت: \lr{RBC}ها تعداد آنتی‌ژن \lr{A} و \lr{B} کمتری دارند؛
\lr{Anti-A, B} به سلول‌های غیر \lr{RBC} که آنتی‌ژن \lr{A} و \lr{B} دارند نیز متصل می‌شود.
\item شایع‌ترین علت هایپربیلی‌روبینمی نیازمند درمان در نوزادان:
\AHmark{همولیز \lr{ABO}}؛ چون همولیز \lr{Rh} به خوبی درمان کنترل‌شده.
\end{itemize}
\subsection*{اریتروبلاستوزیس فتالیس و \lr{Rh}}
\begin{itemize}
\item اریتروبلاستوزیس فتالیس به طور کلاسیک ناشی از \AHmark{ناسازگاری \lr{Rh}} است؛
معمولاً بیماری بارداری دوم به بعد است.
\item آنتی‌ژن‌های سیستم \lr{Rh}: \lr{C, D, E, c, e}؛
\lr{d} خاصیت آنتی‌ژنی ندارد.
\item تست کمبس مثبت می‌شود: \AHmark{\lr{Anti-D Ab}}.
\end{itemize}
\subsection*{اولین کودک تحت تأثیر، مثلاً دومین بارداری}
\begin{itemize}
\item بیماری خفیف: آنمی خفیف + هایپربیلی‌روبینمی.
\item جنین‌های بعدی: بیماری شدید، چون شروع همولیز زودتر رخ می‌دهد.
\item علائم: آنمی جنین، \lr{HF}، افزایش فشار وریدی، انسداد ورید پورت، هایپوآلبومینمی؛
\lr{Fetal Hydrops}.
\end{itemize}

\AHpage{۴}
\subsection*{هیدروپس و شدت همولیز}
\begin{itemize}
\item علائم هیدروپس: آسیت، پلورال و پریکاردیال افیوژن، \lr{Anasarca}.
\item حاشیهٔ \lr{Anasarca}: تجمع زیاد مایع به کل بدن.
\item ریسک مرگ جنین بالاست.
\item بررسی شدت همولیز: \AHmark{میزان بیلی‌روبین مایع آمنیوتیک}.
\end{itemize}
\subsection*{پیشگیری از حساس شدن مادر}
\begin{itemize}
\item \lr{anti-Rh-positive Ig (RhoGAM)}.
\item \AHmark{در بالای هفتهٔ ۲۸ بارداری و ۷۲ ساعت پس از زایمان}.
\item میزان خونی که از جنین طی زایمان وارد جریان خون مادر می‌شود:
\lr{30 mL}؛ \lr{300 \(\mu\)g RhoGAM}.
\item \lr{RhoGAM} صرفاً برای آنتی‌ژن \lr{D} است.
\item گروه‌های خونی دیگری که می‌توانند هیدروپس و اریتروبلاستوز ایجاد کنند:
\lr{Rh C, E, Kell, Duffy}.
\item \AHmark{آلوایمونیزیشن علیه \lr{Kell}}: اریتروپوئز مهار می‌شود؛
عدد رتیک و بیلی‌روبین مایع آمنیوتیک کاهش می‌یابد.
\end{itemize}
\subsection*{علل غیرایمن همولیز در نوزادان}
\begin{itemize}
\item \lr{RBC Enzyme deficiencies}:
\lr{Pyruvate kinase deficiency} و \lr{G6PD deficiency}.
\item اختلالات غشای \lr{RBC}: \lr{HS}؛ فرم \lr{AD} خفیف،
\AHmark{فرم \lr{AR} شدید}؛ کمبود اسپکترین در غشا.
\item هموگلوبینوپاتی‌ها: تالاسمی.
\end{itemize}

\AHpage{۵}
\subsection*{از دست دادن خون}
\begin{itemize}
\item حاد: پارگی آمبلیکال کورد، \lr{fetal-maternal hemorrhage}،
پلاسنتا پرویا، خونریزی داخلی.
\item علائم خونریزی حاد: رنگ‌پریدگی، نبض محیطی ضعیف، شوک.
\item \AHmark{شواهد ارگانومگالی و خون‌سازی خارج \lr{BM} وجود ندارد.}
\item عدد \lr{Hb} ظرف \lr{24 h} افت می‌کند.
\item مزمن: \lr{TTTS}، خونریزی فعال جنین؛ رنگ‌پریدگی واضح،
\AHmark{\lr{HF}}، \AHmark{ارگانومگالی \AHmath{\pm} هیدروپس}.
\item آنمی هایپوکروم میکروسیتیک، کاهش ذخایر آهن.
\end{itemize}
\subsection*{تظاهرات همولیز}
\begin{itemize}
\item همولیز شدید: هیدروپس شدید با \AHmark{آناسارکا}، \lr{HF}، ادم ریوی؛ آسفیکسی.
\item نوزادی که در جنینی همولیز کند: \AHmark{ظرف ۲۴ ساعت اول زردی}؛
رنگ‌پریدگی و \AHmark{ارگانومگالی} دارد.
\item ناسازگاری \lr{ABO}: آنمی خفیف، \AHmark{بدون ارگانومگالی}؛
زردی جزئی \AHmark{ظرف ۲۴–۷۲ ساعت اول}.
\end{itemize}
\subsection*{برای نوزاد با شواهد همولیز چه آزمایش‌هایی؟}
\begin{itemize}
\item \AHmark{\lr{CBC, PBS}}، \AHmark{رتیک، گروه خونی، دایرکت کمبس}.
\item حاشیهٔ دایرکت کمبس: \lr{To determine Ab coated RBCs}.
\item در \lr{PBS}، \AHmark{پلی‌کرومازی و آنیزوسیتوز} قابل انتظار است؛ اگر همولیز ایمنی باشد.
\item پلی‌کرومازی: \lr{RBC} نابالغ زیاد است؛ آنیزوسیتوز: \lr{RBC} با اندازه‌های متفاوت.
\item در ناسازگاری \lr{ABO}، \AHmark{اسفروسیت} در \lr{PBS}؛
حاشیه: \lr{RBC} کروی‌شکل.
\end{itemize}

\AHpage{۶}
\subsection*{همولیز غیرایمنی}
\begin{itemize}
\item شک به همولیز غیرایمنی: کمبس منفی و عدم تفاوت گروه خونی؛
\AHmark{بررسی‌های آنزیمی}.
\item به طور معمول در همولیز مزمن: آنمی شدید + \lr{HF}؛ در خونریزی حاد: شوک.
\item برای بررسی همولیز غیرایمنی: \AHmark{سونوگرافی مغز، کبد، طحال، آدرنال}.
\end{itemize}
\subsection*{درمان آنمی نوزاد}
\begin{itemize}
\item اگر علامت دارد: تزریق خون \lr{Packed}.
\item اگر همولیز وجود داشته باشد، \lr{Pack} باید هم برای مادر و هم نوزاد کراس‌مچ شود.
\item \AHmark{\lr{10--15 mL/kg}} خون به نوزاد می‌دهیم.
\item خون باید اشعه‌دیده باشد و \lr{CMV} منفی؛ \lr{GVHD}.
\item \lr{HIV, HBV, HCV}، سیفلیس.
\item اگر نوزاد در رحم خون گرفته باشد، باید \AHmark{اریتروپوئتین} دریافت کند؛ چون مهار می‌شود.
\end{itemize}
\subsection*{اپروچ به آنمی در نوزاد}
\begin{itemize}
\item \AHmark{عدد رتیک پایین}، ۰–۲٪: \AHmark{\lr{BMA}}.
\item رتیک نرمال یا بالا: \AHmath{>5\text{--}10\%}.
\item کمبس مثبت: ایزوایمونیزیشن؛ کمبس منفی: \AHmark{\lr{PBS}}.
\end{itemize}

\AHpage{۷}
\subsection*{\lr{PBS} در اپروچ به آنمی}
\begin{itemize}
\item دیسمورفولوژی خاص: \lr{HS} و غیره.
\item هایپوکروم میکروسیتیک: آلفا تالاسمی، \lr{TTTS}، گاما تالاسمی، خونریزی.
\item مورفولوژی نرمال، بدون زردی: \AHmark{خونریزی حاد، \lr{DIC}، سپسیس}.
\item مورفولوژی نرمال، با زردی: عفونت‌ها، آنزیموپاتی‌ها، دیابت ملیتوس، داروها
(\AHmark{والپروات})، \AHmark{لوکمی مادرزادی}، \AHmark{گالاکتوزمی}.
\end{itemize}
\subsection*{هایپربیلی‌روبینمی}
\begin{itemize}
\item بیلی‌روبین ناشی از کاتابولیسم هموگلوبین: غیرکونژوگه، محلول در چربی، غیرمستقیم؛
\AHmark{سمی برای \lr{CNS}}؛ به آلبومین متصل می‌شود.
\item \AHmark{یوروبیلینوژن} توسط باکتری‌ها از بیلی‌روبین مستقیم تولید شده،
از \AHmark{ادرار و مدفوع} دفع می‌شود.
\item بازجذب بیلی‌روبین در روده را محدود می‌کند.
\item تأخیر در دفع مکونیوم (که حاوی بیلی‌روبین است): بازجذب بیلی‌روبین.
\item \AHmark{بیلی‌روبین مستقیم از جفت عبور نمی‌کند.}
\end{itemize}
\subsection*{اتیولوژی هایپربیلی‌روبینمی غیرمستقیم}
\begin{itemize}
\item[۱)] زردی فیزیولوژیک: \lr{Diagnosis of exclusion}؛ همولیز، عفونت، بیماری‌های متابولیک.
\item بیلی‌روبین غیرمستقیم در \AHmark{روز سوم} تولد حداکثر \AHmark{\lr{12 mg/dL}}.
\item در نوزادان نارس، در روز پنجم حداکثر \lr{15 mg/dL}.
\item عدد ۱۲ روز سوم: حداکثر در نوزاد شیرمادر \AHmark{۱۵–۱۷}.
\end{itemize}

\AHpage{۸}
\subsection*{موارد پاتولوژیک}
\begin{itemize}
\item \AHmark{اگر تظاهر بالینی زردی در روز نخست تولد وجود داشته باشد: قطعاً پاتولوژیک.}
\item سایر موارد پاتولوژیک: افزایش بیلی‌روبین بیش از \lr{0.5 mg/dL/hour}.
\item حداکثر میزان بیلی‌روبین بیش از \AHmark{\lr{13 mg/dL}} در نوزاد ترم.
\item \lr{Direct bilirubin fraction \(>1.5\) mg/dL}.
\item \AHmark{وجود هپاتواسپلنومگالی یا آنمی}.
\end{itemize}
\subsection*{ادامهٔ اتیولوژی غیرمستقیم}
\begin{itemize}
\item[۲)] سندرم کریگلر–نجار: \AHmark{وراثت \lr{AR}}، کمبود آنزیم کبدی؛
\AHmark{افزایش بیلی‌روبین غیرمستقیم}.
\item تیپ \lr{I}: بسیار شدید، منجر به \AHmark{کرنیکتروس}.
\item تیپ \lr{II}: خفیف؛ \AHmark{فنوباربیتال} منجر به افزایش فعالیت آنزیم.
\item[۳)] بیماری ژیلبرت: هایپربیلی‌روبینمی غیرمستقیم خفیف؛
اگر به طور همزمان همولیز باشد، زردی پیدا می‌کند.
\item[۴)] زردی \lr{Breast milk}: در \AHmark{دو هفتهٔ اول} تولد، هایپربیلی‌روبینمی غیرمستقیم.
\item سطح \lr{Bil} به ندرت بالای \AHmark{\lr{20 mg/dL}} می‌رسد.
\item قطع شیردهی برای ۲ روز: افت شدید \lr{Bil}؛ سپس با شروع مجدد شیردهی عود نمی‌کند.
\item[۵)] زردی روز اول: همیشه پاتولوژیک، مستقیم، غیرمستقیم؛
معمولاً دلیل همولیز یا عفونت؛ خونریزی‌های داخلی.
\item اگر ناشی از عفونت باشد، \AHmark{معمولاً مستقیم} است.
\end{itemize}
\subsection*{مشاهدهٔ زردی و بررسی بیشتر}
\begin{itemize}
\item نوزاد: بیلی‌روبین حداقل \lr{5--10 mg/dL} است؛ در بزرگسالان حداقل \lr{2--3 mg/dL} است.
\item[\AHmath{\ast}] \AHmark{سطح \lr{Bili} بالای \lr{5 mg/dL} در روز نخست و بالای
\lr{13 mg/dL} بعد از آن}: \AHmark{اندیکاسیون بررسی بیشتر}.
\end{itemize}

\AHpage{۹}
\subsection*{بررسی بیشتر زردی}
\begin{itemize}
\item[\AHmath{\ast}] \AHmark{بیلی مستقیم و غیرمستقیم، گروه خون، تست کمبس، \lr{CBC, PBS}، رتیک}.
\item زردی که \AHmark{پس از هفتهٔ دوم} تولد شروع شود: \AHmark{قطعاً پاتولوژیک}؛ معمولاً مستقیم است.
\end{itemize}
\subsection*{اتیولوژی غیرمستقیم}
\begin{itemize}
\item همولیز دارد؛ شایع: عدم تطابق گروه خون (\lr{ABO, Rh, Kell})، دارو و عفونت.
\item ناشایع: آنزیموپاتی‌ها و اختلالات غشا، هموگلوبینوپاتی
\AHunclear{نماد کنار این واژه: \lr{B0} یا بتا}؛
\lr{G6PD, HS}، تالاسمی.
\item همولیز ندارد؛ شایع: زردی فیزیولوژیک، خونریزی، زردی شیر مادر، پلی‌سیتمی، نوزاد مادر \lr{DM}.
\item ناشایع: کریگلر–نجار، ژیلبرت، \AHmark{هایپوتیروئیدی}، \AHmark{\lr{ITP}}.
\end{itemize}
\subsection*{اتیولوژی مستقیم}
\begin{itemize}
\item \AHmark{بیلی‌روبین مستقیم بیش از \lr{2 mg/dL} یا بیش از ۲۰٪ توتال}.
\item \AHmark{همیشه پاتولوژیک است؛ نوروتوکسیک نیست.}
\item[\AHmath{\ast}] \AHmark{بررسی‌ها}: آنزیم‌های کبدی:
\lr{AST, ALT, ALK-P, \(\gamma\)-GT}.
\item کشت ادرار باکتری و ویروس.
\item بررسی‌های متابولیک.
\item سونوگرافی کبد.
\item \AHmark{تست کلر عرق}.
\item \AHmark{ادرار تیره و مدفوع روشن + زردی پس از هفتهٔ دوم}: \AHmark{بیلیاری آترزی}.
\item فتوتراپی و تعویض خون در هایپربیلی مستقیم جایگاهی ندارد.
\item عفونت‌های \AHmark{\lr{TORCH, UTI}}: هایپربیلی \AHmark{مستقیم} ایجاد می‌کنند.
\end{itemize}

\AHpage{۱۰}
\subsection*{کرنیکتروس، انسفالوپاتی بیلی‌روبین}
\begin{itemize}
\item اثر شدید در بازال گانگلیا.
\item در نوزاد ترم: بیلی‌روبین بالای \AHmark{\lr{20--25 mg/dL}}.
\item مگر در شرایط سپسیس، دهیدریشن، گرسنگی، هایپوکسی، هایپوترمی و غیره.
\item اگر وزن هنگام تولد زیر \lr{1 kg} باشد، حتی بیلی زیر ۱۰ نیز خطر دارد.
\end{itemize}
\subsection*{تظاهرات کرنیکتروس}
\begin{itemize}
\item تظاهرات اولیه: لتارژی، هایپوتونی، تحریک‌پذیری، \AHmark{رفلکس مورو کاهش‌یافته}، پورفیدینگ.
\item \AHmark{گریه و جیغ \lr{high pitched}}.
\item پس از \AHmark{روز ۴ام تولد}؛ واژهٔ خط‌خوردهٔ پیش از «پس از»:
\AHunclear{ظاهر؟}.
\item علائم بعدی: فونتانل بالج، خونریزی پولمونری، تب، هایپرتونی، \AHmark{تشنج}.
\item \lr{Paralysis, upward gaze, Opisthotonic posturing}.
\item عوارض: مرگ، \AHmark{ناشنوایی}، \lr{CP, MR, enamel dysplasia}، تغییر رنگ دندان‌ها.
\end{itemize}
\subsection*{درمان هایپربیلی غیرمستقیم}
\begin{itemize}
\item فتوتراپی: قبل از آنکه سطح \lr{Bili} به حد کرنیکتروس برسد باید شروع شود.
\item در نوزاد ترم در روز اول و دوم اگر سطح بیلی غیرمستقیم به بالای \AHmark{\lr{14--16 mg/dL}}.
\item تعویض خون: \AHmark{اگر عدد \lr{Bili} غیرمستقیم به ۲۰ برسد}؛ اگر وزن نوزاد بالای \lr{2 kg} باشد.
\item اگر علامت‌دار نباشد و علت زردی فیزیولوژیک یا شیر مادر باشد:
تا بیلی ۲۵ می‌شود صبر کرد.
\item اگر وزن \lr{1.5 kg}: سطح ۱۵؛ وزن \lr{1 kg}: سطح ۱۰.
\item ۸–۹ ساعت پس از تعویض خون، سطح \lr{Bili} به دلیل ادامهٔ همولیز بالا می‌رود.
\end{itemize}
\subsection*{پلی‌سیتمی، سندرم هایپرویسکوزیته}
\begin{itemize}
\item هماتوکریت بیش از ۶۵٪.
\end{itemize}

\AHpage{۱۱}
\subsection*{گروه در معرض خطر پلی‌سیتمی}
\begin{itemize}
\item نوزادان ترم و پست‌ترم، \lr{SGA}.
\item نوزاد از مادر \lr{DM}.
\item \AHmark{تأخیر در کلمپ بند ناف}.
\item هایپرتیروئیدی نوزادی.
\item تریزومی ۱۳ و ۱۸.
\item \lr{TTTS}، دریافت‌کننده.
\item \AHmark{\lr{Beckwith-Wiedemann Syndrome}}.
\end{itemize}
\subsection*{علائم، تصویر و درمان پلی‌سیتمی}
\begin{itemize}
\item علائم پلی‌سیتمی: \AHmark{آکروسیانوز، تشنج}، آپنه، تحریک‌پذیری،
\AHmark{هایپربیلی‌روبینمی} (\AHmark{ایسکمی کبد}).
\item فشار خون ریوی مداوم نوزاد: \AHmark{\lr{PPHN}}.
\item اگر \lr{CXR} بگیریم: \AHmark{کاردیومگالی، پلورال افیوژن، ادم بینابینی}.
\item درمان: \lr{partial exchange transfusion}؛ جایگزینی بخشی از خون با \lr{NS}.
\end{itemize}
\subsection*{اختلالات انعقادی}
\begin{itemize}
\item \AHmark{فاکتورهای انعقادی از جفت عبور نمی‌کنند.}
\item سطح فاکتورهای وابسته به \lr{vit K} در نوزاد پایین است: \AHmark{۲، ۷، ۹، ۱۰}.
\item در نوزاد، میزان فیبرینوژن (\lr{I})، \lr{V, VIII, XIII} و پلاکت در رنج \lr{NL} بزرگسال است.
حاشیهٔ اعداد رومی: ۱، ۵، ۸، ۱۳.
\item فاکتورهای تماسی (\lr{XI, XII}) در نوزاد پایین است؛ حاشیه: ۱۱، ۱۲.
\item \AHmark{\lr{PTT} در نوزاد طولانی است.}
\item وابسته به: ۱، ۲، ۵، ۸، ۹، ۱۰، ۱۲.
\item \AHmark{\lr{PTT}} نوزاد ترم: \AHmark{\lr{35--50 s}}؛ نوزاد پره‌ترم: \AHmark{\lr{50--80 s}}.
بازهٔ خط‌خوردهٔ پیشین در سطر \lr{PTT} نوزاد ترم:
\AHunclear{\lr{45--65 s}}؛ با ۳۵–۵۰ جایگزین شده است.
\item چه چیزهایی \lr{PTT} را طولانی می‌کند؟ هپارین، \AHmark{\lr{DIC}}، \AHmark{هموفیلی}، کمبود شدید \lr{vit K}.
\end{itemize}

\AHpage{۱۲}
\subsection*{\lr{PT}، \lr{BT} و ویتامین \lr{K}}
\begin{itemize}
\item \lr{PT} در نوزاد ترم به صورت جزئی طولانی است: \lr{13--20 s}.
\item وابسته به فاکتورهای \AHmark{۱، ۲، ۵، ۸، ۱۰}؛ برای \AHmark{\lr{vit K} حساس‌تر است}.
% Factor 8 is transcribed as written in pic12; not silently changed to factor 7.
\item \lr{PT} در چه مواردی طولانی می‌شود؟ کمبود \lr{vit K}، آسیب کبدی، \lr{DIC}.
\item \AHmark{\lr{BT} در نوزاد طبیعی است}؛ نشان‌دهندهٔ تعداد و عملکرد پلاکت.
\item \lr{Breast milk} منبع ضعیف از نظر \lr{vit K} است؛
بالاترین ریسک خونریزی طی روزهای \AHmark{۳–۷ تولد}.
\item تجویز وارفارین، \lr{AB} (ایزونیازید، ریفامپین)، فنوباربیتال، فنی‌توئین
برای مادر باردار: خونریزی روز اول در نوزاد.
عبارت خط‌خوردهٔ پیش از فنوباربیتال: \AHunclear{برست ملک؟}.
\item برای همهٔ نوزادان: \AHmark{تک‌دوز \lr{IM}، \lr{1 mg vit K}}.
\item درمان خونریزی ناشی از کمبود \lr{vit K}: \lr{1 mg vit K IV}؛ \lr{FFP}.
\end{itemize}
\subsection*{ترومبوسایتوپنی و تشخیص افتراقی}
\begin{itemize}
\item یک تشخیص افتراقی مهم در \AHmark{ترومبوسایتوپنی نوزاد \lr{ill}}:
\lr{Kasabach-Meritt Syndrome}؛ \AHmark{گیر افتادن پلاکت در همانژیوم}.
\item \lr{Bacterial infection, NEC, PPHN}، اسیدمی، پره‌اکلامپسی، آسفیکسی، اندوکاردیت ترومبوتیک.
\item اتیولوژی \lr{DIC} در نوزاد: هایپوکسی، هایپوتنشن، آسفیکسی، سپسیس،
ویروس‌ها (به‌خصوص \lr{HSV})، \lr{NEC}، مرگ قل دیگر.
\item هموفیلی: خونریزی جلدی خفیف؛ علائم بیماری سیستمیک ندارد.
\item[\AHmath{\ast}] \AHmark{ترومبوسایتوپنی} در یک \AHmark{\lr{Well child}} ممکن است یک سندرم باشد؛
\AHmark{سندرم آنمی فانکونی}: \AHmark{آپلازی \lr{thumb}}.
\end{itemize}

\AHpage{۱۳}
\subsection*{ادامهٔ سندرم‌ها}
\begin{itemize}
\item \AHmark{سندرم آپلازی رادیال–ترومبوسیتوپنی}: \lr{thumb} وجود دارد.
\item \AHmark{سندرم \lr{Wiskott-Aldrich}}.
\end{itemize}
\subsection*{شایع‌ترین علت ترومبوسایتوپنی در نوزاد \lr{Well}}
\begin{itemize}
\item \AHmark{\lr{anti plt Ab}}:
\lr{Transient isoimmune thrombocytopenia}؛ \lr{Transient neonatal thrombocytopenia}.
\item اگر پلاکت زیر ۲۰٬۰۰۰ بیاید یا خونریزی جدی داشته باشد:
\begin{latin}
Transfusion of ABO compatible, RhD compatible, PA-1a-neg
and PA-5b-neg maternal plt.
\end{latin}
\item حتی اگر درمان نشود، ظرف ۱ ماه \lr{plt} به رنج \lr{NL} برمی‌گردد.
\item زایمان \lr{C/S}؛ ریسک \lr{ICH} را کم می‌کند.
\end{itemize}
\subsection*{ترومبوسیتوپنی نوزادی در نوزاد با مادر \lr{ITP}}
\begin{itemize}
\item ناشی از عبور آنتی‌بادی \lr{IgG} از جفت.
\item درمان مادر با کورتیکواستروئید؛ زایمان \lr{C/S}.
\item درمان نوزاد: پردنیزولون، \lr{IVIG}.
\item در صورت عدم درمان هم ظرف یک ماه از بین می‌رود.
\end{itemize}
\subsection*{بررسی خونریزی نوزاد}
\begin{itemize}
\item \AHmark{هر نوزادی با خونریزی}: \AHmark{\lr{Plt, PBS, PT, PTT}}.
\item \AHmark{ترومبوسایتوپنی ایزوله در یک نوزاد \lr{Well}}:
\AHmark{آلوایمیون ترومبوسایتوپنی ایمنی}.
\item \lr{DIC}: \AHmark{\lr{PT, PTT} به شدت بالا} (در حد دقیقه)؛ ترومبوسایتوپنی.
\item \lr{PBS}: \AHmark{بورسل، فرگمنتهٔ \lr{RBC}}
(میکروآنژیوپاتیک همولیتیک آنمی).
\item کمبود \lr{vit K}: \AHmath{PT\uparrow\uparrow\uparrow}، \AHmath{PTT\uparrow}.
\item هموفیلی (کمبود ۸ و ۹): \AHmark{صرفاً افزایش \lr{PTT}}.
\end{itemize}

\clearpage
\section*{پیوست: تصاویر کامل منبع}
متن اصلی، حاشیه‌ها، فلش‌ها، خط‌خوردگی‌ها، هایلایت‌ها و سربرگ‌های هر ۱۳ برگه در اصل تصاویر زیر محفوظ‌اند.
\IfFileExists{images/pic01.jpg}{\def\AHimagebase{images/}}{%
  \def\AHimagebase{pediatrics/ped-neonatology/anemia & hyperbilirubinemia/images/}}
\newcommand{\AHoriginal}[2]{%
  \clearpage\subsection*{اصل تصویر #1: \lr{\texttt{pic#2.jpg}}}%
  \begin{center}%
  \includegraphics[width=\linewidth,height=0.87\textheight,keepaspectratio]{\AHimagebase pic#2.jpg}%
  \end{center}}
\AHoriginal{۱}{01}
\AHoriginal{۲}{02}
\AHoriginal{۳}{03}
\AHoriginal{۴}{04}
\AHoriginal{۵}{05}
\AHoriginal{۶}{06}
\AHoriginal{۷}{07}
\AHoriginal{۸}{08}
\AHoriginal{۹}{09}
\AHoriginal{۱۰}{10}
\AHoriginal{۱۱}{11}
\AHoriginal{۱۲}{12}
\AHoriginal{۱۳}{13}
\AHfinish
